%\documentclass[DiscStateSpaceToolkit.tex]{subfiles} 
%\begin{document}
 
Aiyagari (1994) is somewhat similar to Hugget (1993), the later has a pure-exchange economy while
the first is a production economy, but the models otherwise have much in common. Both are interested
in the question of how idiosyncratic shocks, incomplete markets, and borrowing constraints interact and
their possible effects on interest rates.

In this model agents labour supply is exogenously determined (the idiosyncratic shocks act via varying wage 
earings). Agents aim to maximize their consumption by saving capital. They are unable to borrow, but can
accumulate the capital as precautionary savings. In the aggregate these individual savings determine the aggregate
capital level (and their exogenous labour supplies the aggregate labour supply), which in turn, via the assumption
of a Cobb-Douglas production function and competitive markets (so factors are paid marginal products) determine 
the interest rate and wage.


Thus, the agents problem is to chose consumption $c$ and assets $k$ to maximize
\begin{align*}
 E[\sum_{t=0}^{\infty} \beta^t \frac{c^{1-\mu}-1}{1-\mu}]
\end{align*}
(note that for $\mu=1$ this is log-utility) subject to a budget constraint and no-borrowing constraint
\begin{align*}
 c_t+k_{t+1}=w l_t + (1+r)k_t \\
 c_t \geq 0, \text{ } a_t \geq -b
\end{align*}
From the assumption of a Cobb-Douglas production function $f=K^{\alpha}L^{1-\alpha}$ and perfect competition, and that capital
depreciates at rate $\delta$, we have that
\begin{align}
 r=\alpha K^{\alpha-1} L^{1-\alpha} -\delta \label{Aiyagari:eqn:r} \\
 w=(1-\alpha) K^{\alpha} L^{-\alpha} \label{Aiyagari:eqn:w}
\end{align}
The labour endowment shocks, $l_t$, are assumed to follow an AR(1) process in logs, namely
\begin{align*}
 log(l_t)=\rho log(l_{t-1})+\sigma(1-\rho^2)^{1/2} \epsilon_t        \qquad \qquad \epsilon_t \sim N(0,1)
\end{align*}

There are currently two markets that must clear, capital and labour (we need to get the interest rate and wage to be the relevant marginal products), but 
before coding the model notice that in fact we don't need to try and solve for both interest rate and wage since there is a unique mapping between them. This
comes about because of the use of a Cobb-Douglas production function; combining equations \ref{Aiyagari:eqn:r} \& \ref{Aiyagari:eqn:w} we get that in equilbrium
$w=(1-\alpha)((r+\delta)/\alpha)^{-\alpha /(1-\alpha)}$. Thus we only need a grid for interest rates (not wages) and only need to look at market clearing in
the capital market (to check eqn \ref{Aiyagari:eqn:r} holds). Another thing to notice is that since labour is exogenous and there are a continuum of agents
of weight one we will always have that $L=E[l]=E[exp(log(l))]$ (recall that the AR(1) process on $l$ is in logs). Thus we don't need to go and calculate $L$ 
again for every change in prices (interest rates and wages). The borrowing constraint is easily imposed by simply having the assets grid start from $-b$.

Aiyagari defines as the 'natural borrowing limit', the maximum amount that could be paid back at the equilibrium interest rate if the agent were forever 
stuck in the low wage state (in the models notation; natural borrowing limit is $b_n=wl_1/r$, where $l_1$ is the smallest state of $l_t$). Any borrowing 
limit stricter than this is called an 'ad-hoc' borrowing limit. In particular the case of 'no borrowing' ($b=0$) used in the computation of the model
constitutes an ad-hoc borrowing limit.

The code implementing the model is following, loosely speaking it recreates individual elements of Tables I \& II of that paper. Aiyagari also looks 
at a couple of other results; recreating them using the objects created by the code is fairly easy but not done here.


\lstinputlisting[language=Matlab]{./Examples/MatlabCodes/Aiyagari1994.m}

%\end{document}